\subsection{Algorithm Design}
\label{sec:algorithm_design}

Two algorithmic decisions shape what a learner may do next: when a card is
scheduled for review, and when a stage of a career pathway opens. The first is
the hybrid SM-2 model, whose full specification --- the derivation of $Q$ from
objective response signals, the calibration weight, the interval rules, and the
learner-facing retention metrics built on them --- is given in
Appendix~\ref{appendix:algorithm}. It is placed there because this iteration
does not evaluate it: a scheduling policy's effect on retention can only be
established over a cohort observed across months, which is outside the scope of
the present work. The pathway rules below are specified here because they are
not a research claim but a piece of academic bookkeeping the system must get
exactly right, and because every progress figure the product reports is derived
from them.

\subsubsection{Career Pathway Progress and Completion}
\label{sec:pathway_completion}
The scheduling model governs movement \textit{within} a course. A career pathway is the structure above it: an ordered sequence of stages, each bundling required courses with a set of electives from which the stage asks for a stated number. Progress along a pathway, and the decision that a student has finished one, are therefore computed over stages rather than over cards.

\paragraph{Stage Satisfaction and the Elective Quota}
Let $\mathcal{S}$ denote the stages of the career-path version the student is pinned to. For a stage $s$, let $R_s$ be its required courses, $O_s$ its electives, and $m_s \in \mathbb{N}_0$ the number of electives the stage requires — its quota. A course $c$ is \textit{satisfied}, written $\sigma(c) = 1$, when the student holds a completion award for it; partial lesson progress does not satisfy a course.

The work a stage credits and the work it demands are then

\begin{equation}
d_s = \left|\left\{ c \in R_s \mid \sigma(c) = 1 \right\}\right|
     + \min\left(\left|\left\{ c \in O_s \mid \sigma(c) = 1 \right\}\right|,\; m_s\right),
\qquad
t_s = \left|R_s\right| + m_s.
\end{equation}

The $\min$ is the whole of the quota rule: electives satisfied beyond $m_s$ are credited at $m_s$ and no further, so a stage offering five electives and asking for one is finished by the first of them. This is a deliberate departure from treating every attached course as owed. An elective set exists to offer a choice, and a design under which choosing one obliges the student to the other four would make that choice meaningless. Because $d_s \leq t_s$ holds by construction, a stage can never be credited beyond what it demands.

\paragraph{Pathway Progress}
Pathway progress aggregates the stages, excluding those that demand nothing:

\begin{equation}
P = \frac{\sum_{s \in \mathcal{S},\; t_s > 0} d_s}{\sum_{s \in \mathcal{S},\; t_s > 0} t_s} \times 100.
\end{equation}

A stage with $t_s = 0$ — no required courses and no quota — has nothing outstanding and nothing to measure. Leaving it in the denominator would hold a finished pathway permanently below $100$, so it is excluded from both sums rather than counted as incomplete.

\paragraph{Completion Is a Predicate}
Whether a student has finished a pathway is decided from the stages directly,

\begin{equation}
\text{complete} \iff \mathcal{S} \neq \emptyset \;\wedge\; d_s = t_s \quad \forall s \in \mathcal{S},
\end{equation}

rather than from a presentation metric. Since $d_s \leq t_s$ for every stage,
the predicate and the aggregate agree exactly: $P = 100$ if and only if every
stage is complete. The two are one rule under two presentations, which is what
prevents a pathway reading as finished on screen while the records built on
completion remain open.

An empty $\mathcal{S}$ is excluded deliberately. A pathway with no stages has not been shown to be finished; it has merely offered nothing to examine, and what completion writes is a durable academic record. Publication requires at least one stage per pathway and at least one course per stage, so the case does not arise for any pathway a student can be enrolled onto.

\paragraph{Stage Unlocking}
Completion says what a student has finished; unlocking says what they may open. The two are computed from the same evaluation but answer different questions, and a pathway needs both: a stage may be finished without ever having been reachable, and the design turns on recognising that as a coherent state rather than a contradiction.

Each stage carries an unlock policy. Writing $U(s)$ for "stage $s$ is unlocked", $C(s)$ for "stage $s$ is complete" and $R(s)$ for "every required course of $s$ is satisfied", the rule for the stage at position $i$ is

\begin{equation}
U(s_i) =
\begin{cases}
\text{true} & i = 1 \text{ or policy} = \texttt{always},\\
C(s_{i-1}) \wedge U(s_{i-1}) & \text{policy} = \texttt{after\_previous},\\
R(s_{i-1}) \wedge U(s_{i-1}) & \text{policy} = \texttt{after\_previous\_required},\\
\text{false} & \text{otherwise.}
\end{cases}
\end{equation}

The first stage is unconditionally open whatever its stored policy says: a pathway whose opening stage were locked could not be started by anyone. An unrecognised policy resolves to false rather than true, so a malformed rule withholds content instead of releasing it. The distinction between the two gated policies is the elective quota: \texttt{after\_previous} waits for the previous stage to be complete, quota included, while \texttt{after\_previous\_required} releases the next stage once its required work is done and lets the electives run on.

\paragraph{Why Unlocking Chains}
The conjunction with $U(s_{i-1})$ is the part that is easy to omit and costly to omit. Satisfaction is a property of a student and a course, with no pathway in it, so a stage's courses can be satisfied by work done somewhere else entirely. Where a programme permits several pathways at once, a second pathway whose course set mirrors one stage of the first will complete that stage in the first the moment it is joined --- correctly, because the work was genuinely done, and it counts toward that pathway and is recorded as permanently earned.

It must not, however, open what follows it. Without the conjunction, a stage the student had never been permitted to enter would serve as the prerequisite pass for the stage after it: a curriculum whose first stage is a hard requirement could be entered at its third stage without the first ever being touched, by taking a second pathway that mirrors the second stage. Requiring the previous stage to have been reachable as well as finished closes that route, and does so without weakening the credit: the mirrored stage remains complete, remains counted, and remains permanently recorded.

What this admits is a stage that is complete and locked at the same time. That pair is not a contradiction but the precise description of the situation --- the work is done and will never have to be done again, while the pathway's own sequence has not yet reached it. An \texttt{always} stage is unaffected throughout, being unlocked by its own policy; it propagates the chain rather than interrupting it, so an unfinished stage of pure electives can never shut the pathway behind it.

Unlocking governs presentation and access rather than eligibility to progress. A locked stage blocks a student from starting its courses only where its author marked the stage as hard-enforced; the softer settings record the lock, warn the student, and let them through, because an author who chose to advise rather than to block should not be silently overruled.

\paragraph{One Evaluation, Every Surface}
The same evaluation answers the question wherever it is asked: on the learner's own pathway view, on the manager's roster of every student following a pathway, and in the frozen snapshot written when a student leaves one. None of these recomputes progress for itself.

This is a stronger requirement than it appears, because the three surfaces have different shapes and invite different shortcuts. A roster reports many students at once, and a single aggregate query is the obvious way to produce it: each student's percentage, and the count of courses behind it, derived in one statement from the mean lesson completion of every attached course. Such a figure is reasonable in isolation and wrong in company. A course counts as satisfied when the student has been awarded its completion, which is not the same condition as its lessons averaging $100$, so a roster built that way can report a student at $100\%$ against a course count that says otherwise, in the same row.

The correction is not a better query but one fewer. The roster's query now returns identity and avatar columns only; the percentage, the number of satisfied courses and the denominator are all read from the stage evaluation, which is the sole definition any of them has. A second derivation of a value that already has one authority does not add a check — it adds a disagreement that nothing is watching for.

Evaluating per student is what this costs, since satisfaction is per student by nature. The stage structure, however, is a property of the pathway version and not of the reader, so it is resolved once for the whole roster rather than re-read for each row.

\paragraph{Programme Completion}
Completion propagates upward but never sideways. Finishing a pathway completes the attempt that held it, and the programme enrolment behind that attempt completes only once none of its attempts remain active — so a student holding two pathways completes the programme on the second, not the first. An attempt ended by an approved change or drop is closed rather than active, and therefore does not hold the enrolment open. Each attempt is judged against the pathway version pinned to it, so a revision published mid-enrolment cannot alter what an already-enrolled student must finish.

\paragraph{The evaluation in one pass}
Figure~\ref{fig:stage_evaluation_flow} in Appendix~\ref{appendix:algorithm}
traces the rules above as the evaluator applies them: the ordered stage walk,
the quota arithmetic, the latch that preserves a completion already earned, and
the forward dependency by which each stage's unlock reads the stage before it.
